% from ega4-8-fr.tex
\oldpage[IV-3]{35}
Nous en verrons de nombreux exemples par la suite et nous nous bornerons ici
à donner quelques résultats de ce type.
\end{env}

\begin{proposition}[8.9.3]
\label{IV.8.9.3-fr}
Soient $A$ un anneau, $M$ un $A$-module de présentation finie. Alors tout
$A$-endomorphisme surjectif de $M$ est bijectif.
\end{proposition}

\begin{proof}
En effet, considérons $A$ comme limite inductive de ses sous-$\mathbf Z$-algèbres
de type fini. Il résulte de (8.5.2.6) qu'il existe une de ces sous-algèbres
$A_0$ et un $A_0$-module $M_0$ de présentation finie tels que $M$ soit
$A$-isomorphe à $M_0\otimes_{A_0}A$~; en outre, si $f:M\to M$ est un
$A$-endomorphisme surjectif, on peut supposer (8.5.2, (i)) qu'il existe un
$A_0$-endomorphisme $f_0:M_0\to M_0$ tel que $f=f_0\otimes1_A$~; enfin
(8.5.7) on peut supposer que $f_0$ est surjectif. Mais comme $A_0$ est
noethérien et $M_0$ un $A_0$-module de type fini, $M_0$ est un $A_0$-module
noethérien, donc (Bourbaki, \emph{Alg.}, chap. VIII, \S~2, n\textsuperscript{o}~2,
lemme 3) $f_0$ est bijectif, et il en est par suite de même de $f$.
\end{proof}

\begin{proposition}[8.9.4]
\label{IV.8.9.4-fr}
(«~théorème de platitude générique~»). --- Soient $Y$ un préschéma intègre,
$u:X\to Y$ un morphisme de type fini et localement de présentation finie,
$\mathcal F$ un $\mathcal O_X$-Module quasi-cohérent de présentation finie. Il
existe alors un ouvert non vide $U$ de $Y$ tel que
$\mathcal F|u^{-1}(U)$ soit plat sur $U$.
\end{proposition}

\begin{proof}
Le raisonnement du début de (6.9.1) ramène à prouver le

\begin{lemma}[8.9.4.1]
\label{IV.8.9.4.1-fr}
Soient $A$ un anneau intègre, $B$ une $A$-algèbre de présentation finie,
$M$ un $B$-module de présentation finie. Alors il existe un $f\neq0$ dans
$A$ tel que $M_f$ soit un $A_f$-module libre.
\end{lemma}

En effet, d'après (8.9.1) il y a une sous-$\mathbf Z$-algèbre de type fini
$A_0$ de $A$, une $A_0$-algèbre de type fini $B_0$ et un $B_0$-module de type
fini $M_0$ tels que $B$ soit $A$-isomorphe à $B_0\otimes_{A_0}A$ et que $M$
soit $B$-isomorphe à $M_0\otimes_{A_0}A$~; d'après (6.9.2), il existe
$f_0\neq0$ dans $A_0$ tel que $(M_0)_{f_0}$ soit un
$(A_0)_{f_0}$-module libre. Comme
$M_{f_0}=(M_0)_{f_0}\otimes_{A_0}A$ et
$A_{f_0}=(A_0)_{f_0}\otimes_{A_0}A$, $M_{f_0}$ est un $A_{f_0}$-module libre.
\end{proof}

\begin{corollary}[8.9.5]
\label{IV.8.9.5-fr}
Soient $S$ un préschéma quasi-compact et quasi-séparé, $u:X\to S$ un
morphisme de présentation finie, $\mathcal F$ un $\mathcal O_X$-Module
quasi-cohérent de présentation finie. Il existe alors une partition
$(S_i)_{1\leq i\leq n}$ de $S$ en un nombre fini d'ensembles localement fermés
dans $S$ telle que, pour $1\leq i\leq n$, il existe un sous-préschéma $S'_i$
de $S$, ayant pour espace sous-jacent $S_i$, de présentation finie sur $S$,
tel que si l'on pose $X_i=X\times_SS'_i$, le $\mathcal O_{X_i}$-Module
$\mathcal F_i:=\mathcal F\otimes_{\mathcal O_S}\mathcal O_{S'_i}$ soit plat
sur $S'_i$.
\end{corollary}

\begin{proof}
Considérons un recouvrement fini $(U_j)_{1\leq j\leq m}$ de $S$ par des
ouverts affines et définissons par récurrence $T_1=U_1$,
$T_j=U_j-\bigcup_{h<j}(U_j\cap U_h)$ pour $j\geq2$~; chaque $T_j$ est fermé
dans l'ouvert affine $U_j$, et les $T_j$ sont deux à deux sans point commun~;
en outre les $U_j\cap U_h$ sont quasi-compacts puisque $S$ est quasi-séparé,
et par suite ($S$ étant aussi quasi-compact), sont constructibles dans $S$
(I, 1.8.1), donc il en est de même des $T_j$. Il suffira évidemment de
démontrer le corollaire pour un sous-préschéma convenable $T'_j$ de $S$ ayant
$T_j$ pour espace sous-jacent, de présentation finie sur $S$ et pour le
morphisme et le Module déduits respectivement de $u$ et $\mathcal F$ par le
changement de base $T'_j\to S$. Notons pour cela que si l'on prend sur $U_j$
la structure de préschéma induite par celle de $S$, l'immersion ouverte
$U_j\to S$ est quasi-compacte puisque $S$ est quasi-séparé (I, 1.2.7), donc
elle est de
